\section{Fazit}
Im Rahmen dieser Arbeit konnten wir zeigen, dass sich die Erzeugung eines 3D Pythagoras-Baums gut als datenparalleles Problem
modellieren lässt, da zwischen den Quadern keine Abhängigkeiten bestehen und somit kaum Synchronisation nötig ist. Entscheidend 
für die erfolgreiche Parallelisierung war die Umstellung von einer rekursiven auf eine iterative, ebenenweise Berechnung, weil sich 
nur so die Arbeit sinnvoll auf CPU-Kerne bzw. GPU-Threads verteilen lässt.\\
Unsere Laufzeitanalyse zeigte, dass der Overhead paralleler Ansätze für kleine Baumtiefen (insbesondere bei OpenCL durch 
Speichertransfer und Kernel-Setup) sehr groß ist, während bei größeren Tiefen deutliche Zeitgewinne möglich sind. OpenMP erwies 
sich als überaus nützlich, während OpenCL zwar bei großen Tiefen Vorteile bringen kann, jedoch stark von 
Hardware/Implementierung abhängt (z. B. PoCL-Effekt durch Wegfall des GPU-Transfers).\\
Der zentrale limitierende Faktor unseres Projekts ist nicht die Rechenleistung, sondern der Speicher. Bereits ab ca. 21 Ebenen
liegt der Bedarf allein für die Quaderdaten im Gigabyte-Bereich, bei 25 Ebenen im zweistelligen Gigabyte-Bereich.
Deshalb waren Bäume mit mehr als ca. 23 Ebenen in der Praxis nicht mehr wirklich generierbar.\\
Als naheliegende Weiterentwicklungen bieten sich insbesondere eine Pipeline-/Paketierung der OpenCL-Transfers, das Reduzieren 
verbleibender kritischer Bereiche sowie eine stärkere Parallelisierung bzw. Beschleunigung des Renderers an.